% AUDIT DE VERACITE 2026-09-20 -- corrections des citations de rapport de jury.
% Bloc 4, encadre jury :
%   - SUPPRIME la citation attribuee au « Rapport jury 2019 » sur le « couple thermique equivalent » :
%     verbatim introuvable dans les 8 rapports disponibles (2018-2025) ; la locution « couple thermique »
%     elle-meme n'apparait dans aucun rapport. Citation fabriquee, aucun remplacement pertinent trouve.
%   - REMPLACE la citation attribuee au « Rapport jury 2024 » (plage de mesure / tolerance du capteur),
%     egalement introuvable, par le verbatim authentique du Rapport_Jury_2020.md ligne 703, qui traite
%     bien de la justification du choix du capteur. Annee de citation corrigee en consequence.
% Aucun autre contenu de la fiche n'a ete modifie.
\expandafter\ifx\csname CAPETINCLUDEMODE\endcsname\relax
\documentclass[11pt,a4paper]{article}
% preamble.tex — packages et commandes communes a toutes les fiches CAPET SII IE
\input{glyphtounicode}
\pdfgentounicode=1
\usepackage{cmap}
\usepackage[T1]{fontenc}
\usepackage[utf8]{inputenc}
\usepackage[french]{babel}
% babel[french] active a \begin{document} des raccourcis de ponctuation (:;!?)
% qui cassent la syntaxe to[...] de circuitikz (erreur "+ or - expected") ;
% on les desactive, ce qui ne change que l'espacement fin avant ces signes.
\AtBeginDocument{\shorthandoff{:;!?}}
\usepackage[a4paper, top=1.6cm, bottom=1.6cm, left=1.4cm, right=1.4cm]{geometry}
\setlength{\headheight}{14.5pt}
\usepackage{amsmath, amssymb, amsfonts}
\usepackage{siunitx}
\sisetup{locale=FR, output-decimal-marker={,}, inter-unit-product=\ensuremath{{}\cdot{}},
         group-minimum-digits=4}
% Unites non fournies par siunitx, declarees ici une fois pour toutes (fiches + exercices).
% Les gardes \ifdefined evitent tout conflit avec une redeclaration locale restee dans un
% fichier : ne pas retirer les gardes, ne pas redeclarer ces unites dans les fichiers.
\DeclareSIUnit{\jour}{j}
\ifdefined\var\else\DeclareSIUnit{\var}{var}\fi   % puissance reactive
\ifdefined\vad\else\DeclareSIUnit{\vad}{VAD}\fi   % puissance deformante
\ifdefined\tr\else\DeclareSIUnit{\tr}{tr}\fi      % tour (tr/min)
\usepackage{array, booktabs, multirow, tabularx}
\usepackage{xcolor}
\usepackage[most, breakable]{tcolorbox}
\usepackage{multicol}
\usepackage{enumitem}
\usepackage{mathtools}
\usepackage{lmodern}
% Refonte visuelle 2026 (direction "moderne coloré") : Lato pour le texte,
% maths sans-serif assorties via newtxsf (charge apres amsmath/amssymb).
\usepackage[default]{lato}
\usepackage{newtxsf}
\usepackage{microtype}
\usepackage{fancyhdr}
\usepackage{tikz}
\usetikzlibrary{positioning,arrows.meta}
\usepackage{circuitikz}
\usepackage{lastpage}

% ── Couleurs par domaine (palette modernisée 2026) ────────────────────
\definecolor{colElec}{RGB}{29, 78, 216}        % Électrotechnique — bleu vif
\definecolor{colElecLight}{RGB}{219, 234, 254} % fond bleu clair
\definecolor{colAuto}{RGB}{21, 128, 61}        % Automatique — vert
\definecolor{colMeca}{RGB}{194, 65, 12}        % Mécanique — orange
\definecolor{colTherm}{RGB}{185, 28, 28}       % Thermique — rouge
\definecolor{colQT}{RGB}{232, 244, 233}        % (hérité, plus utilisé par questiontype)
\definecolor{colJury}{RGB}{255, 248, 225}      % fond jury (crème)
\definecolor{colJuryFrame}{RGB}{202, 138, 4}   % bordure/bandeau jury (ambre)
\definecolor{colGrayLight}{RGB}{245, 245, 247} % fond gris neutre
\definecolor{colSeparator}{RGB}{180, 180, 180} % gris séparateur
\definecolor{colExo}{RGB}{8, 145, 178}          % Banque d'exercices — cyan (distinct des fiches)
% ── Encre et gris de la refonte ───────────────────────────────────────
\definecolor{colInk}{RGB}{23, 32, 42}          % texte courant (quasi-noir)
\definecolor{colSlate}{RGB}{83, 98, 118}       % gris ardoise (méta, en-têtes)
\definecolor{colRule}{RGB}{203, 213, 225}      % filets clairs
\definecolor{colJuryBg}{RGB}{255, 248, 225}    % fond crème jury
\definecolor{colPresence}{RGB}{83, 98, 118}    % bloc 5 (présence dans les sujets)
\AtBeginDocument{\color{colInk}}

% ── Style encadré exercice + corrigé (banque d'exercices) ─────────────
\tcbset{
  exoBox/.style={
    breakable,
    enhanced,
    colback=colExo!8!white,
    colframe=colExo,
    colbacktitle=colExo,
    coltitle=white,
    fonttitle=\bfseries\small,
    left=4pt, right=4pt, top=3pt, bottom=3pt,
    before skip=5pt, after skip=5pt,
  }
}

% ── Bascule enonce / corrige (banque d'exercices) ─────────────────────
% \CAPETCORRIGE est defini en ligne de commande (pdflatex "\def\CAPETCORRIGE{1}\input{...}")
% pour compiler la variante corriges d'un meme fichier source ; par defaut (non defini),
% la compilation produit la variante enonces. \CAPETSHOW{enonce}{corrige} choisit la branche.
%
% TROIS MODES (refonte 2026-09-20, manuel de revision) :
%   - aucun drapeau   -> imprime l'enonce            (variante enonces, INCHANGE)
%   - \CAPETCORRIGE   -> imprime le corrige          (variante corriges, INCHANGE)
%   - \CAPETMANUEL    -> imprime l'enonce ET met le corrige de cote pour l'annexe
%
% Le mode manuel stocke chaque corrige dans une macro (\csgdef), sans ecriture
% fichier ni re-tokenisation : le contenu (circuitikz, accents, %) traverse intact.
% L'annexe les rejoue ensuite dans l'ordre via \CAPETANNEXE.
\usepackage{etoolbox}

\newcounter{CAPETexo}          % numero continu d'exercice dans le manuel
\setcounter{CAPETexo}{0}
\newcounter{CAPETannexe}       % compteur de relecture en annexe

% Titre courant de la notion en cours (pose par \CAPETNOTION), sert d'intertitre annexe
\newcommand{\CAPETnotioncode}{}
\newcommand{\CAPETnotiontitre}{}

% \CAPETNOTION{code}{titre} : ouvre une notion dans le manuel (memorise le libelle,
% et marque le rang du premier exercice de la notion pour l'intertitre d'annexe).
\newcommand{\CAPETNOTION}[2]{%
  \renewcommand{\CAPETnotioncode}{#1}%
  \renewcommand{\CAPETnotiontitre}{#2}%
  \csgdef{CAPETnotionstart\the\numexpr\value{CAPETexo}+1\relax}{#1 -- #2}%
}

% ── \degree robuste ───────────────────────────────────────────────────
% Certains fichiers d'exercices utilisent \degree. Ni le preambule ni siunitx
% ne le definissent, et \providecommand local ne suffit pas : sous hyperref,
% \degree existe deja comme commande factice (\ERROR), donc \providecommand
% s'abstient et l'appel echoue a la compilation du manuel (erreur invisible
% en compilation isolee). On le (re)definit donc inconditionnellement ici,
% apres le chargement des packages, pour tout le corpus.
\AtBeginDocument{\providecommand{\degree}{}\renewcommand{\degree}{\ensuremath{^\circ}}}

% Libelle decoratif : les fichiers d'exercices utilisent aussi \CAPETSHOW dans
% leur \fancyhead et leur bloc de titre, pour afficher "(enonces)" ou "(corriges)".
% Ces appels-la ne sont PAS des exercices : en mode manuel ils ne doivent ni
% consommer un numero, ni alimenter l'annexe. On les reconnait a leur premier
% argument, strictement egal au libelle ci-dessous.
\newcommand{\CAPETLIBELLEENONCES}{(\'enonc\'es)}

\ifdefined\CAPETMANUEL
  % --- Mode manuel : enonce en place, corrige differe vers l'annexe ---
  \newcommand{\CAPETSHOW}[2]{%
    \def\CAPETargun{#1}%
    \ifx\CAPETargun\CAPETLIBELLEENONCES
      #1%            appel decoratif : rendu tel quel, hors numerotation
    \else
      \CAPETSHOWEXO{#1}{#2}%
    \fi
  }
  \newcommand{\CAPETSHOWEXO}[2]{%
    \stepcounter{CAPETexo}%
    \csgdef{CAPETcorrige\theCAPETexo}{#2}%
    \edef\CAPETtmpnum{\theCAPETexo}%
    #1%
    \par\vspace{2pt}%
    {\footnotesize\itshape\color{colExo}%
     \hyperref[capet:cor:\CAPETtmpnum]{Corrige de l'exercice \CAPETtmpnum\ $\rightarrow$ p.~\pageref{capet:cor:\CAPETtmpnum}}%
     \label{capet:eno:\CAPETtmpnum}}%
  }
  % \CAPETANNEXE : rejoue tous les corriges captures, dans l'ordre, avec intertitres
  \newcommand{\CAPETANNEXE}{%
    \setcounter{CAPETannexe}{0}%
    \loop\ifnum\value{CAPETannexe}<\value{CAPETexo}%
      \stepcounter{CAPETannexe}%
      \edef\CAPETcurnum{\theCAPETannexe}%
      \ifcsname CAPETnotionstart\CAPETcurnum\endcsname
        \par\medskip
        {\normalsize\bfseries\color{colExo}\csuse{CAPETnotionstart\CAPETcurnum}}%
        \par\nobreak\smallskip
      \fi
      \begin{tcolorbox}[exoBox, breakable,
        title={Corrige de l'exercice \CAPETcurnum\
               \hfill\normalfont\footnotesize
               \hyperref[capet:eno:\CAPETcurnum]{enonce p.~\pageref{capet:eno:\CAPETcurnum}}}]%
      \label{capet:cor:\CAPETcurnum}%
      \csuse{CAPETcorrige\CAPETcurnum}%
      \end{tcolorbox}%
    \repeat
  }
\else
  % --- Modes historiques, comportement strictement inchange ---
  \newcommand{\CAPETSHOW}[2]{\ifdefined\CAPETCORRIGE #2\else #1\fi}
  \newcommand{\CAPETANNEXE}{}
\fi

% ── Difficulté 1 à 3 (étoiles pleines, amssymb) ───────────────────────
\newcommand{\difficulte}[1]{%
  \ifcase#1\relax
  \or $\bigstar$\;niveau 1/3%
  \or $\bigstar\bigstar$\;niveau 2/3%
  \or $\bigstar\bigstar\bigstar$\;niveau 3/3%
  \fi
}

% ── Style encadré extrait de document technique (banque d'exercices) ──
\tcbset{
  docTech/.style={
    breakable,
    enhanced,
    colback=colGrayLight,
    colframe=colSeparator,
    fonttitle=\bfseries\small\itshape,
    coltitle=black,
    colbacktitle=colGrayLight,
    left=4pt, right=4pt, top=3pt, bottom=3pt,
    before skip=4pt, after skip=4pt,
  }
}

% ── Style tcolorbox de section principale ────────────────────────────
% Refonte 2026 : bandeau titre coloré arrondi + fond très pâle, sans cadre.
\tcbset{
  ficheSection/.style={
    breakable, enhanced jigsaw,
    frame hidden, boxrule=0pt, arc=1.8mm,
    colback=#1!4!white,
    colbacktitle=#1,
    coltitle=white,
    fonttitle=\sffamily\bfseries,
    left=7pt, right=7pt, top=5pt, bottom=5pt,
    toptitle=3pt, bottomtitle=3pt,
    before skip=10pt, after skip=8pt,
  }
}

% ── Style encadré question type ──────────────────────────────────────
% Carte blanche, bandeau sombre, liseré très clair.
\tcbset{
  questiontype/.style={
    breakable, enhanced jigsaw,
    colback=white,
    frame hidden, boxrule=0pt, arc=1.8mm,
    colbacktitle=colElec!78!black,
    coltitle=white,
    fonttitle=\sffamily\bfseries\small,
    left=7pt, right=7pt, top=4pt, bottom=4pt,
    toptitle=2.5pt, bottomtitle=2.5pt,
    before skip=8pt, after skip=8pt,
    borderline={0.6pt}{0pt}{colElec!12!white},
  }
}

% ── Style encadré jury ───────────────────────────────────────────────
% Bandeau ambre, fond crème, sans cadre.
\tcbset{
  jury/.style={
    breakable, enhanced jigsaw,
    colback=colJuryBg,
    frame hidden, boxrule=0pt, arc=1.8mm,
    colbacktitle=colJuryFrame,
    coltitle=white,
    fonttitle=\sffamily\bfseries,
    left=7pt, right=7pt, top=4pt, bottom=4pt,
    toptitle=3pt, bottomtitle=3pt,
    before skip=10pt, after skip=8pt,
  }
}

% ── En-tête / pied de page communs (refonte 2026) ─────────────────────
% #1 code fiche, #2 titre court, #3 couleur domaine
\newcommand{\ficheHeaders}[3]{%
  \fancyhead[L]{\small\sffamily{\bfseries\textcolor{#3}{#1}}\hspace{0.7em}\textcolor{colSlate}{#2}}%
  \fancyhead[R]{\small\sffamily\textcolor{colSlate}{CAPET SII IE\, --- \,Session 2026}}%
  \fancyfoot[C]{\small\sffamily\textcolor{colSlate}{\thepage\,/\,\pageref{LastPage}}}%
  \fancyhead[C]{}\fancyfoot[L]{}\fancyfoot[R]{}%
}

% ── Bloc en-tête de fiche (BLOC 0, refonte 2026) ──────────────────────
% #1 code  #2 titre  #3 couleur domaine  #4 domaine  #5 ligne présence  #6 sous-thèmes
\newcommand{\ficheEntete}[6]{%
  \begin{tcolorbox}[enhanced, breakable,
    frame hidden, boxrule=0pt, arc=2.5mm,
    interior style={left color=#3, right color=#3!75!black},
    left=9pt, right=9pt, top=7pt, bottom=7pt,
    before skip=0pt, after skip=6pt]
    {\color{white}\sffamily\bfseries\LARGE #1\quad #2}\\[4pt]
    {\color{#3!22!white}\sffamily\small \MakeUppercase{#4}\quad\textbar\quad #5}\\[1pt]
    {\color{#3!35!white}\sffamily\footnotesize Sous-thèmes : #6}
  \end{tcolorbox}%
}

% ── Mise en page ─────────────────────────────────────────────────────
\pagestyle{fancy}
\fancyhf{}
\renewcommand{\headrulewidth}{0.5pt}
\renewcommand{\footrulewidth}{0pt}
\AtBeginDocument{%
  \renewcommand{\headrule}{{\color{colRule}\hrule width\headwidth height 0.5pt}}%
}

\setlist[itemize]{leftmargin=1.2em, itemsep=1pt, topsep=2pt, parsep=0pt}
\setlist[enumerate]{leftmargin=1.2em, itemsep=3pt, topsep=2pt, parsep=0pt}

\setlength{\parskip}{2pt}
\setlength{\columnsep}{10pt}

\def\CAPETBEGINFILE{\begin{document}}
\def\CAPETENDFILE{\end{document}}
\else
\def\CAPETBEGINFILE{}
\def\CAPETENDFILE{}
\fi

\ficheHeaders{T05}{Phénomènes thermiques}{colTherm}

\CAPETBEGINFILE

% -----------------------------------------------------------------------
%  BLOC EN-TETE
% -----------------------------------------------------------------------

\ficheEntete{T05}{Phénomènes thermiques}{colTherm}%
  {Thermique / Instrumentation}%
  {\textbf{Présence :} \textbf{7/13} sujets (2013-2026) - thème \emph{fréquent}}%
  {couple thermique équivalent - résistance thermique - capteurs (PT100, thermocouple) - transmetteur 4-20\,mA - vapeur saturée}

% -----------------------------------------------------------------------
%  1. RAPPELS DE COURS
% -----------------------------------------------------------------------

\begin{tcolorbox}[
  ficheSection=colTherm,
  title={1. Rappels de cours},
]

% -- 1.1 Couple thermique equivalent ----------------------------------
\textbf{1.1 Couple thermique équivalent sur un cycle}

\smallskip
Un moteur soumis à un cycle de couples $C_{m1}, C_{m2}, \ldots$ de durées $\Delta t_1, \Delta t_2, \ldots$ (période totale $T$) est
caractérisé par son couple thermique équivalent :
\[
  C_{th} = \sqrt{\frac{\sum_{i} C_{mi}^2 \cdot \Delta t_i}{T}}
\]
Le courant thermique équivalent $I_{th}$ s'obtient de même ($I_{th}^2 = \sum I_i^2\Delta t_i/T$).
$C_{th}$ doit être inférieur au couple nominal du moteur pour garantir l'échauffement admissible.

\medskip
\hrule\vspace{4pt}

% -- 1.2 Resistance thermique -----------------------------------------
\textbf{1.2 Résistance thermique et modèle thermique}

\smallskip
\begin{center}
\renewcommand{\arraystretch}{1.45}
\begin{tabular}{llll}
\toprule
\textbf{Grandeur} & \textbf{Expression} & \textbf{Unité} & \textbf{Analogie électrique} \\
\midrule
Flux thermique (conduction) & $\Phi = \lambda S\,\Delta T/e$ & W & courant \\[4pt]
Résistance thermique (conduction) & $R_{th} = e/(\lambda S)$ & K/W & résistance \\[4pt]
Résistance thermique (convection) & $R_{th,conv} = 1/(hS)$ & K/W & résistance \\[4pt]
Élévation de température & $\Delta T = R_{th} \cdot P$ & K & loi d'Ohm \\[4pt]
Température jonction & $T_j = T_{amb} + R_{th,ja} \cdot P_{dissipee}$ & $^\circ\mathrm{C}$ & \\[4pt]
Constante de temps & $\tau_{th} = R_{th} \cdot C_{th,cap}$ & s & $\tau = RC$ \\
\bottomrule
\end{tabular}
\end{center}

\smallskip
\textbf{Circuit thermique équivalent (analogie électrique) :}
\begin{center}
\begin{circuitikz}[european, scale=0.9]
  \draw (0,0) node[left]{$T_j$} to[R, l=$R_{th}$, i>=$\Phi$] (3.4,0) node[right]{$T_{amb}$};
\end{circuitikz}
\end{center}
Le flux thermique $\Phi$ joue le rôle du courant, l'écart de température $\Delta T = T_j - T_{amb}$ celui de la tension : $\Delta T = R_{th} \cdot \Phi$, exactement comme la loi d'Ohm.

\medskip
\hrule\vspace{4pt}

% -- 1.3 Capteurs de temperature --------------------------------------
\textbf{1.3 Capteurs de température}

\smallskip
\begin{center}
\renewcommand{\arraystretch}{1.35}
\begin{tabular}{lllll}
\toprule
\textbf{Capteur} & \textbf{Principe} & \textbf{Plage} & \textbf{Sortie} & \textbf{Précision} \\
\midrule
PT100 & Résistif ($R$ croît avec $T$) & $-200$ à $+850^\circ\mathrm{C}$ & $\Omega$ & $\pm 0{,}3^\circ\mathrm{C}$ \\[3pt]
Thermocouple J & Seebeck (Fe/Constantan) & $-30$ à $+700^\circ\mathrm{C}$ & $\mu$V/K & $\pm1{,}5^\circ\mathrm{C}$ cl.1 \\[3pt]
Thermocouple K & Seebeck (NiCr/Ni) & $-200$ à $+1260^\circ\mathrm{C}$ & $\mu$V/K & $\pm 2{,}2^\circ\mathrm{C}$ cl.1 \\[3pt]
DS1820 & Numérique, bus 1-Wire & $-55$ à $+125^\circ\mathrm{C}$ & bits & $\pm0{,}5^\circ\mathrm{C}$ \\
\bottomrule
\end{tabular}
\end{center}

\smallskip
PT100 : $R(T) = 100\,(1 + \alpha T)$ avec $\alpha = 3{,}85\times10^{-3}\,\text{K}^{-1}$.
Thermocouple : $e_{AB} = k_S(T_1 - T_2)$ (effet Seebeck). \textbf{Toujours compenser la jonction froide} si $T_{ref} \neq 0^\circ\mathrm{C}$.

\medskip
\hrule\vspace{4pt}

% -- 1.4 Transmetteur 4-20 mA -----------------------------------------
\textbf{1.4 Transmetteur 4-20 mA}

\smallskip
Relation linéaire entre grandeur physique $T$ et courant $I_{out}$ :
\[
  I_{out} = 4 + 16 \cdot \frac{T - T_{min}}{T_{max} - T_{min}} \quad [\text{mA}]
\]
$4\,\text{mA} \leftrightarrow T_{min}$ ; $20\,\text{mA} \leftrightarrow T_{max}$ ; $\Delta I = 16\,\text{mA}$ sur toute la plage.

\end{tcolorbox}

% -----------------------------------------------------------------------
%  2. FORMULAIRE SYNTHETIQUE
% -----------------------------------------------------------------------

\begin{tcolorbox}[
  ficheSection=colTherm,
  title={2. Formulaire synthétique},
]
\small
\begin{multicols}{2}
\renewcommand{\arraystretch}{1.45}
\begin{tabular}{@{}p{4.2cm}ll@{}}
\toprule
\textbf{Formule} & \textbf{Usage} & \textbf{U.} \\
\midrule
\multicolumn{3}{@{}l}{\footnotesize\textit{Thermique moteur}} \\
$C_{th}^2 = \sum C_{mi}^2\Delta t_i\,/\,T$ & Couple therm. & N$\cdot$m \\
$I_{th}^2 = \sum I_i^2\Delta t_i\,/\,T$ & Courant therm. & A \\
\midrule
\multicolumn{3}{@{}l}{\footnotesize\textit{Résistance thermique}} \\
$R_{th} = e/(\lambda S)$ & Conduction & K/W \\
$R_{th} = 1/(hS)$ & Convection & K/W \\
$\Delta T = R_{th} \cdot P$ & Échauffement & K \\
$T_j = T_{amb} + R_{th,ja} \cdot P$ & Jonction & $^\circ\mathrm{C}$ \\
\bottomrule
\end{tabular}

\columnbreak

\renewcommand{\arraystretch}{1.45}
\begin{tabular}{@{}p{4.2cm}ll@{}}
\toprule
\textbf{Formule} & \textbf{Usage} & \textbf{U.} \\
\midrule
\multicolumn{3}{@{}l}{\footnotesize\textit{Capteurs}} \\
$R(T) = 100(1+\alpha T)$ & PT100 & $\Omega$ \\
$e = k_S(T_1 - T_2)$ & Thermocouple & V \\
\midrule
\multicolumn{3}{@{}l}{\footnotesize\textit{Transmetteur 4-20\,mA}} \\
$I_{out} = 4 + 16\cdot(T-T_{min})/\Delta T_m$ & Conversion & mA \\
$\Delta T_m = T_{max} - T_{min}$ & Plage capteur & $^\circ\mathrm{C}$ \\
\midrule
\multicolumn{3}{@{}l}{\footnotesize\textit{Effet Joule}} \\
$P_J = R\,I^2 = U\,I$ & Source chaleur & W \\
$\tau_{th} = R_{th} \cdot C_{th,cap}$ & Cste de temps & s \\
\bottomrule
\end{tabular}
\end{multicols}
\normalsize
\end{tcolorbox}

% -----------------------------------------------------------------------
%  3. QUESTIONS TYPES
% -----------------------------------------------------------------------

\begin{tcolorbox}[
  ficheSection=colTherm,
  title={3. Questions types - extraites des sujets réels},
]

% -- QT1 -------------------------------------------------------------

\begin{tcolorbox}[questiontype,
  colframe=colTherm!55!black, colbacktitle=colTherm!55!black,
  title={QT1 - Couple thermique équivalent sur un cycle \hfill\normalfont\footnotesize E2 sujet 2019}
]
\textbf{Énoncé.}
Un moteur subit un cycle de période $T = 2048\,\text{s}$.
Les phases 1 et 2 sont caractérisées par $F_{e1} = 14{,}9\,\text{kN}$, $\Delta t_1 = 42\,\text{s}$ et $F_{e2} = 9{,}2\,\text{kN}$, $\Delta t_2 = 51\,\text{s}$.
Pour les phases 3 à 18 : $\sum C_{mi}^2\Delta t_i = 558\times10^6\,\text{N}^2\text{m}^2\text{s}$.
Calculer le couple thermique équivalent $C_{th}$.

\tcbline

\textbf{Corrigé.}

\textbf{Étape 1 — Contribution des phases 1 et 2.}
(Les couples $C_{m1}$, $C_{m2}$ sont calculés à partir des forces et du bras de levier fournis en DT.)
\[
  \sum_\text{phases 1 et 2} C_{mi}^2\Delta t_i = C_{m1}^2 \cdot 42 + C_{m2}^2 \cdot 51
\]

\textbf{Étape 2 — Couple thermique total.}
\[
  C_{th} = \sqrt{\frac{C_{m1}^2\cdot42 + C_{m2}^2\cdot51 + 558\times10^6}{2048}}
\]

\textbf{Conclusion.} Le couple thermique équivalent obtenu doit être comparé au couple nominal du moteur retenu ; s'il lui est inférieur, l'échauffement reste admissible sur le cycle. \textit{Erreur fréquente : oublier de prendre la racine carrée à la fin (on calcule $C_{th}^2$ et non $C_{th}$).}
\end{tcolorbox}

\vspace{4pt}

% -- QT2 -------------------------------------------------------------

\begin{tcolorbox}[questiontype,
  colframe=colTherm!55!black, colbacktitle=colTherm!55!black,
  title={QT2 - Transmetteur 4-20\,mA - calcul du courant de sortie \hfill\normalfont\footnotesize E2 disciplinaire 2024 (partie E)}
]
\textbf{Énoncé.}
Un transmetteur de température a la plage : $T_{min} = -20^\circ\mathrm{C} \to 4\,\text{mA}$, $T_{max} = +200^\circ\mathrm{C} \to 20\,\text{mA}$.
Calculer $I_{out}$ pour $T = 120^\circ\mathrm{C}$.

\tcbline

\textbf{Corrigé.}

\textbf{Étape 1 — Application de la relation linéaire.}
\[
  I_{out} = 4 + 16 \cdot \frac{T - T_{min}}{T_{max} - T_{min}}
          = 4 + 16 \cdot \frac{120 - (-20)}{200 - (-20)}
          = 4 + 16 \cdot \frac{140}{220}
\]

\textbf{Étape 2 — Résultat.}
\[
  I_{out} = 4 + 10{,}2 \approx \mathbf{14{,}2\,\text{mA}}
\]

\textbf{Conclusion.} $14{,}2\,\text{mA}$ est bien compris entre 4 et 20\,mA, cohérent avec $T = 120^\circ\mathrm{C}$ dans la plage $[-20;200]^\circ\mathrm{C}$ : la mesure est exploitable telle quelle par la chaîne d'acquisition.\checkmark
\end{tcolorbox}

\vspace{4pt}

% -- QT3 -------------------------------------------------------------

\begin{tcolorbox}[questiontype,
  colframe=colTherm!55!black, colbacktitle=colTherm!55!black,
  title={QT3 - Surveillance par PT100 - algorithme et interpolation \hfill\normalfont\footnotesize E2 disciplinaire 2022 (cogénération)}
]
\textbf{Énoncé.}
Trois sondes PT100 ($T_1$, $T_2$, $T_3$) surveillent une canalisation vapeur à $P_s = 12\,\text{bars}$.
Lire la température de vaporisation dans la table et écrire l'algorithme de surveillance.

\tcbline

\textbf{Corrigé.}

\textbf{Étape 1 — Lecture de table.}
À $P_s = 12\,\text{bars}$ : $T_{vap} = 187{,}96^\circ\mathrm{C}$ (interpolation linéaire si nécessaire).

\textbf{Étape 2 — Algorithme de surveillance.}
\begin{center}
\renewcommand{\arraystretch}{1.2}
\begin{tabular}{l}
\textbf{SI} $(T_1 - T_2 < 12^\circ\mathrm{C})$ \textbf{ET} $(T_3 < 160^\circ\mathrm{C})$ \textbf{ALORS}\\
\quad état $\leftarrow$ Fonctionnement correct\\
\textbf{SINON}\\
\quad état $\leftarrow$ Alarme / Défaut\\
\textbf{FIN SI}
\end{tabular}
\end{center}

\textbf{Conclusion.} $(T_1 - T_2 < 12^\circ\mathrm{C})$ détecte un retour d'eau trop froid (condensats) ; $T_3 < 160^\circ\mathrm{C}$ surveille la vapeur en fin de circuit : ces deux seuils couvrent les deux défauts thermiques les plus probables de l'installation.
\end{tcolorbox}

\end{tcolorbox}

% -----------------------------------------------------------------------
%  4. ATTENTES DU JURY
% -----------------------------------------------------------------------

\begin{tcolorbox}[
  jury,
  title={4. Attentes du jury - Erreurs récurrentes et points de vigilance},
]

\begin{enumerate}[label=\textbf{\arabic*.}, leftmargin=*, itemsep=5pt]

  \item \textbf{Oublier la racine carrée dans $C_{th}$.}
  La formule donne $C_{th}^2$ : il faut impérativement extraire $\sqrt{\cdot}$ à la fin.
  S'arrêter au carré est l'erreur la plus signalée par le jury 2019.

  \item \textbf{Transmetteur 4-20\,mA : $\Delta I = 16\,\text{mA}$, pas 20.}
  Le courant varie de 4 à 20\,mA, soit une plage de $16\,\text{mA}$.
  Diviser par 20 au lieu de 16 est une erreur fréquente.

  \item \textbf{Interpolation linéaire dans les tables thermodynamiques.}
  Si la pression exacte n'est pas tabulée, interpoler entre les deux valeurs encadrantes.
  Ne pas prendre la valeur la plus proche sans calcul.

  \item \textbf{Thermocouple : compensation de jonction froide.}
  La tension fournie est $e = k_S(T_{mesure} - T_{ref})$. Si $T_{ref} \neq 0^\circ\mathrm{C}$, appliquer la correction thermique avant de conclure.

  \item \textbf{PT100 vs thermocouple : ne pas confondre les principes.}
  PT100 : résistif (pont de Wheatstone) - grande précision, plage limitée.
  Thermocouple : effet Seebeck - plage étendue, sensibilité moindre.

  \item \textbf{Effet Peltier $\neq$ effet Seebeck.}
  Seebeck : $\Delta T \to e$ (capteur, thermocouple). Peltier : $I \to \Delta T$ (module thermoélectrique). L'effet exploité dans un thermocouple est l'effet Seebeck.

  \item \textbf{Préciser la plage et la tolérance du capteur choisi.}
  Le jury apprécie de vérifier que la grandeur à mesurer est dans la plage du capteur et de préciser la classe de tolérance.

\end{enumerate}

\medskip
\begin{tcolorbox}[
  colback=colJuryFrame!12!colJury,
  colframe=colJuryFrame,
  left=4pt, right=4pt, top=2pt, bottom=2pt,
  fontupper=\small\itshape
]
Rapport jury 2020 : «~Le choix du capteur, dans cette partie, est également très rarement justifié. Le jury rappelle que toute conclusion doit être chiffrée en lien avec une grandeur physique.~»
\end{tcolorbox}

\end{tcolorbox}

% -----------------------------------------------------------------------
%  PIED - PRESENCE DANS LES SUJETS
% -----------------------------------------------------------------------

\begin{tcolorbox}[
  ficheSection=colPresence,
  title={Présence dans les sujets analysés (2013-2026)},
]
\begin{center}
\footnotesize\renewcommand{\arraystretch}{1.15}
\begin{tabular}{lccccccccccccc}
\toprule
\textbf{Année} & 13 & 14 & 15 & 16 & 17 & 18 & 19 & 20 & 21 & 22 & 23 & 24 & 25 \\
\midrule
E1 (commune)   & - & - & - & - & - & - & - & - & - & $\checkmark$ & - & $\checkmark$ & - \\
E2 discipl. IE & - & - & - & $\checkmark$ & - & $\checkmark$ & $\checkmark$ & $\checkmark$ & - & $\checkmark$ & - & $\checkmark$ & $\checkmark$ \\
\bottomrule
\end{tabular}
\end{center}
\smallskip
\footnotesize\textbf{Sous-thèmes fréquents :}
couple thermique équivalent (5/13) - capteurs PT100/thermocouple (7/13) - transmetteur 4-20\,mA (4/13) - tables vapeur saturée (3/13) - résistance thermique composants (3/13)

\medskip
\textbf{Ordres de grandeur (repères rapides).}
\begin{itemize}[leftmargin=1.3em]
  \item Réseau BT FR : tension phase-neutre typique \SI{230}{\volt}, tension phase-phase typique \SI{400}{\volt}, fréquence \SI{50}{\hertz}.
  \item Courants d'étude fréquents en sujet : de quelques \si{\ampere} à quelques dizaines de \si{\ampere}.
  \item En conversion d'énergie, un rendement attendu reste généralement entre 0{,}8 et 0{,}98 selon la technologie.
\end{itemize}

\medskip
\textbf{Synthèse express (à refaire sans support).}
\begin{enumerate}[leftmargin=1.4em]
  \item Poser les hypothèses du modèle et les conventions de signe.
  \item Écrire 2 à 4 relations clés avec unités.
  \item Réaliser un calcul numérique complet et vérifier l'ordre de grandeur.
  \item Justifier le résultat par un argument technique de niveau CAPET.
\end{enumerate}

\end{tcolorbox}

\CAPETENDFILE
